\documentclass[10pt,a4paper]{amsart}
\usepackage[margin=19mm]{geometry}
\usepackage{amsmath,amssymb,mathtools}
\usepackage[scr=rsfs,cal=euler]{mathalfa}
\usepackage{xfrac}
\usepackage[unicode,hidelinks,bookmarks=false]{hyperref}
\newcommand{\ie}{i.e.\ }
\newcommand{\cf}{cf.\ }
\newcommand{\resp}{respectively\ }
\newcommand{\Subsection}{\subsection}
\providecommand{\nobreakdash}{\nobreak\hbox{-}}
\setlength{\emergencystretch}{2em}
\multlinegap0pt
\usepackage{bm}
\usepackage{bbm}
\usepackage{dsfont}
\usepackage{xspace}
\usepackage{cancel}
\usepackage{mathtools}
\usepackage{cleveref}
\usepackage{amsthm}
\makeatletter
\@ifundefined{theorem}{\newtheorem{theorem}{Theorem}}{}
\@ifundefined{cor}{\newtheorem{cor}[theorem]{Corollary}}{}
\@ifundefined{lemma}{\newtheorem{lemma}[theorem]{Lemma}}{}
\@ifundefined{proposition}{\newtheorem{proposition}[theorem]{Proposition}}{}
\makeatother
\DeclareMathOperator{\Ext}{Ext}
\newcommand{\cE}{{\mathcal E}}
\newcommand{\mC}{{\mathbb C}}
\newcommand{\mF}{{\mathbb F}}
\newcommand{\mM}{{\mathbb M}}
\newcommand{\mR}{{\mathbb R}}
\title[Splittings of truncated motivic Brown--Peterson cooperations]{Splittings of truncated motivic Brown--Peterson cooperations algebras}
\author{Jackson Morris, Sarah Petersen,, Elizabeth Tatum}
\date{Source: arXiv:2509.19542v2 (CC BY 4.0)}
\begin{document}
\maketitle

\noindent\emph{Source and licence.} Splittings of truncated motivic Brown--Peterson cooperations algebras. Version arXiv:2509.19542v2 (CC BY 4.0), distributed under the Creative Commons Attribution 4.0 International licence, as recorded in the arXiv licence metadata for this exact version and shown on the abstract page https://arxiv.org/abs/2509.19542v2. Full rights notice: licenses/arxiv-2509.19542-CC-BY-4.0.txt.\par\smallskip

\noindent\textit{Editorial scope.} Counted unit: the Proposition whose source label is \texttt{prop:lightningExtCase1} in the article (source0.tex lines 1780--1782, source label \texttt{prop:lightningExtCase1}), delivered together with the article's own complete original proof of it (source0.tex lines 1784--1829, one \texttt{proof} environment of 46 lines). No mathematics is written, repaired or re-derived: statement and proof are the article's own text line for line. Required context retained uncounted: lightning flash ses (lines 864--867); lem:lightning flash ext 1 (lines 1373--1379); cor: Ext L\_0 L\_m (lines 1442--1446); prop:ExtLightningflashKleqM (lines 1660--1674). Every definition, display and label of those lines is retained unchanged. Omitted from this package and not redistributed: 1--863, 868--1372, 1380--1441, 1447--1659, 1675--1779, 1783--1783, 1830--2189. Nothing inside a retained excerpt is omitted or altered: every retained body is delivered exactly as the article's lines are written, so no omission receipt of any kind accompanies this package. Citations: the retained excerpts carry no citation command at all, so no bibliography entry of the article is transcribed and no reference is redistributed. Reference adaptation: none was needed and none was made; every reference inside the delivered unit resolves inside it, so no unresolved reference or citation is produced. Printed slips kept literal: no printed word, symbol, hypothesis or number of the counted statement or of its proof was altered. Layout and class: the article's own class is \texttt{amsart} with packages including geometry, breakurl, color, tikz-cd, amsmath, amssymb, graphicx, mathrsfs, mathabx, xy, adjustbox, ulem, spectralsequences, tcolorbox; the wrapper is the lane's pinned \texttt{amsart} editorial wrapper at 10pt on A4, which restates the theorem-like environments the retained text needs and adds the article's own macro definitions that the retained text needs (\texttt{\textbackslash Ext}, \texttt{\textbackslash cE}, \texttt{\textbackslash mC}, \texttt{\textbackslash mF}, \texttt{\textbackslash mM}, \texttt{\textbackslash mR}), with extra wrapper packages (bm, bbm, dsfont, xspace, cancel, mathtools, cleveref). The delivered numbering is the unit's own. Assets: no retained span contains a figure environment, a graphics-inclusion command, a PDF-page inclusion, a table or a TikZ picture; no image file is copied into this package, the package declares no asset dependency and no image file is redistributed with it. Licence: version arXiv:2509.19542v2 of this article is distributed under the Creative Commons Attribution 4.0 International licence, as recorded in the arXiv licence metadata for this exact version and shown on the abstract page https://arxiv.org/abs/2509.19542v2. A full rights notice is delivered at licenses/arxiv-2509.19542-CC-BY-4.0.txt. Characters: every retained excerpt is pure ASCII, so the delivered engine, XeLaTeX with its default Latin Modern text font, reports no missing character.
\begin{equation}
\label{lightning flash ses}
0 \to \Sigma^{2(p-1),p-1} L_p(k-1) \to L_p (k) \to \left( \cE(1)_p // \cE(0)_p \right)^\vee \to 0.
\end{equation}

\begin{lemma}
\label{lem:lightning flash ext 1}
    Let $L_p(k)$ denote the $k^{th}$ lightning flash module. Then there is an isomorphism of $\mathbb{M}_p$-modules and $\mathcal{E}(1)_p^\vee$-comodules:
    \[\textup{Ext}^{s,f,w}_{\mathcal{E}(1)_p^\vee}(L_p(k)) \cong \bigoplus_{i=0}^{k-1}\textup{Ext}^{s,f,w}_{\mathcal{E}(0)_p^\vee}(\mathbb{M}_p)\{x_i\} \oplus \textup{Ext}^{s,f,w}_{\mathcal{E}(1)_p^\vee}(\mathbb{M}_p) \{x_k\},\]
    where $|x_i|=(2i(p-1), 0, i(p-1))$ and where \[v_1x_i = v_0x_{i+1}\]
    for all $0 \leq i \leq k$.
\end{lemma}

\begin{cor}\label{cor: Ext L_0 L_m}
    If $F = \mathbb{C}$ at any prime, $F = \mathbb{R}$ and $p$ is odd, or $F = \mF_q$ and there is a trivial Bockstein action, then 
    \[ \Ext_{\mathcal{E}(1)^\vee_p}^{s,f,w}(\mathbb{M}_p^F,L_p(m)) \cong \Ext_{{E}(1)^\vee_p}^{s,f}(\mathbb{F}_p,L_p^{cl}(m)) \otimes \mM_{p}^{F} ,\]
    where $L_p^{cl}(m)$ denotes the classical lightning flash module.
\end{cor}

\begin{proposition}
\label{prop:ExtLightningflashKleqM}
    Suppose that $k \leq m$. Then 
    \[\textup{Ext}^{s,f,w}_{\mathcal{E}(1)^\vee_p}(L_p(k), L_p(m)) \cong \textup{Ext}_{\mathcal{E}(1)^\vee_p}^{s,f,w}(\mathbb{M}_p, L_p(m-k)) \oplus B,\]
    where $B$ consists of $v_0$ and $v_1$-torsion concentrated in filtration $f=0$ and negative stems.
    To be precise, $B$ consists of:
    \begin{enumerate}
        \item A direct sum of $\mathbb{M}_p^{\mathbb{C}}$ in filtration 0 and negative stems congruent to 1 modulo $2(p-1)$, when $F=\mathbb{C}$ and $p$ is any prime;
        \item A direct sum of $\mathbb{M}_p^\mathbb{R}$ in filtration 0 and negative stems congruent to 1 modulo $2(p-1)$, whenever $F=\mathbb{R}$ and $p>2$;
        \item A direct sum of $\mathbb{F}_2[\rho, \tau^4]$ in filtration 0 with generator in negative stems congruent to 1 modulo $2$, whenever $F=\mathbb{R}$ and $p=2$;
        \item A direct sum of $\mathbb{M}_p^{\mathbb{F}_q}[u]/(u^2)$ in filtration 0 with generator in negative stems congruent to 1 modulo $2(p-1)$, whenever $F=\mathbb{F}_q$ has a trivial Bockstein action and $p$ is any prime;
        \item A direct sum of $\mathbb{F}_2[\tau^2]$,  $\mathbb{F}_2[\tau^2, \rho]/(\rho^2)$, and $\mathbb{F}_2[\tau^2, \rho\tau]/((\rho\tau)^2)$ in filtration 0 with generator in negative stems congruent to 1 modulo 2, whenever $F=\mathbb{F}_q$ has a nontrivial Bockstein action and $p=2$;
        \item A direct sum of $\mathbb{F}_p[\zeta^p]$, $\mathbb{F}_p[\zeta^p, \gamma]/(\gamma^2)$, and $\mathbb{F}_p[\zeta^p, \gamma\zeta^j]/((\gamma\zeta^j)^2)$ for all $1 \leq j \leq p-1$ in filtration 0 with generator in negative stems congruent to 1 modulo $2(p-1)$, whenever $F=\mathbb{F}_q$ has a nontrivial Bockstein action and $p>2$.
    \end{enumerate}
\end{proposition}

\begin{proposition}\label{prop:lightningExtCase1}
    Suppose $F = \mC$ or $F = \mF_q$ and $p$ is any prime where $\textup{char}(\mathbb{F}_q) \equiv 1 \,(p^2)$, or $F = \mR$ and $p$ is odd. Then \[\Ext_{\cE(1)_p^\vee}^{s,f,w}(L_p(k), L_p(m)) \cong \Ext_{E(1)_p^\vee}^{s,f}(L^{cl}_p(k), L^{cl}_p(m)) \otimes \mM_{p}^{F} .\]
\end{proposition}

\begin{proof}
    First, note that the case of $k \leq m$ is an immediate consequence of \cref{prop:ExtLightningflashKleqM} combined with the description of  $\text{Ext}_{\mathcal{E}(1)^\vee_p}^{s,f,w}(\mM_p^F, L_p(m))$ for $k \le m$ found in \cref{cor: Ext L_0 L_m}. Thus, we must only prove the case of $k>m$. 
    % Recall that $\Ext_{E(1)_*}^{s,f}(L_p^{cl}(k), L_p^{cl}(m)) $ can be described as follows. When $k \le m$, 
    % \[
    % \Ext_{E(1)_*}^{s,f}(L_p^{cl} (k), L_p^{cl} (m)) \cong \mF_{2}[v_{0}, v_{1}]\{x_{0}, x_{1}, \ldots x_{m-k}|\ v_{1}x_{i} = v_{0}x_{i+1} \},
    % \] 
    % where $|x_{i}| = (2i, 0)$. When $k \ge m$, 
    % \[
    % \Ext_{E(1)_{*}}^{s,f}(L_p^{cl} (k), L_p^{cl} (m)) \cong \mF_{p}[v_{0}, v_{1}]\{x_0\} \oplus \mF_{p}[v_{0}, v_{1}]\left\{ y_{0}, y_{1}, \ldots, y_{k-m-1} \Bigg\vert \begin{array}{l}
    %      v_{1}y_{i} = v_{0}y_{i+1},\\  v_{0}y_{0} = 0,\\
    %      v_{1}y_{k-m} = 0 
    % \end{array}\right\},
    % \]
    % where $|x| = (0, k-m)$ and $|y_{i}| = (-1 - 2(k-m-i), 0)$.

Fix an $m \ge 0$, and start with the base case $k=m$. By \cref{prop:lightningExtCase1},
\[
\textup{Ext}_{\mathcal{E}(1)^\vee_p}^{s,f,w}(L_p(m), L_p(m)) \cong \textup{Ext}_{{E}(1)^\vee_p}^{s,f}(L^{cl}_p(m), L^{cl}_p(m)) \otimes \mM_{p}^{F}.
\]

Suppose that for all $k' < k$, $\Ext_{\cE(1)_p^\vee}^{s,f,w}(L_p^F (k'), L_p^F (m)) \cong \Ext_{E(1)_p^\vee}^{s,f}(L_p^{cl}(k'), L_p^{cl} (m)) \otimes \mM_{p}^{F}$. 
As in the previous proposition, the short exact sequence of lightning flash modules from \Cref{lightning flash ses} induces a long exact sequence of the form
\begin{align}
    \label{eq:LESLightningFlashExtKgeqM}
    \begin{split}
        \cdots \to \text{Ext}^{s,f,w}_{\mathcal{E}(1)^\vee_p}(L_p (k), L_p (m)) \to \text{Ext}^{s,f,w}_{\mathcal{E}(1)^\vee_p}(\Sigma^{2(p-1), p-1}L_p (k-1), L_p (m)) \\
        \xrightarrow{d}\text{Ext}^{s-1,f+1,w}_{\mathcal{E}(1)^\vee_p}((\mathcal{E}(1)_p//\mathcal{E}(0)_p)^\vee, L_p(m)) \to \cdots
    \end{split}
\end{align}

Recall that 
\begin{align*}
    {\Ext}^{s,f,w}_{\mathcal{E}(1)^\vee_p}((\mathcal{E}(1)//\mathcal{E}(0))_p^\vee, L_p(m)) & \cong {\Ext}^{s,f,w}_{\mathcal{E}(0)^\vee_p}(\Sigma^{-2(p-1), -(p-1)}\mM_p^F, L_p(m)) \\
    & \cong \mM_p^F[v_0]\{y\},
\end{align*} where $|y| = (-2(p-1),0, -(p-1))$. Consider the generator 
\[
\Sigma^{-2(p-1), -(p-1)}x \in \text{Ext}^{-2(p-1),0,0}_{\mathcal{E}(1)^\vee_p}(\Sigma^{2(p-1), p-1}L_p(k-1), L_p(m))
\]
in degree $(-2(p-1),0,0)$. Note that the differential $d$ preserves motivic weight, so either $d(x) = v_0^{m-k }y$ or $d(x) = 0$. Comparison with  $\Ext^{s,f,w}_{\cE(1)_p^\vee} \left(L_p (m+1), L_p (m+1)\right)$, which we computed in \cref{prop:ExtLightningflashKleqM}, implies that $d(x)$ must be nonzero. Specifically, suppose towards a contradiction that $d(\Sigma^{-2(p-1), -(p-1)}x) = 0$. Then we will have an infinite $v_{0}$-tower in stem $s=-(2p-1)$ of $\Ext^{s,f,w}_{\mathcal{E}(1)^\vee_p}\left(L_p^F(k), L_p^F(m)\right)$ for $k=m+1$ and $k=m$. Combining this with the long exact sequence used in the inductive computation of \cref{lem:lightning flash ext 1} would imply $\Ext^{s,f,w}_{\mathcal{E}(1)^\vee_p}\left(L_p(m+1), L_p(m+1)\right)$ must also have an infinite $v_{0}$-tower in odd stem. But we already have already proven in \cref{lem:lightning flash ext 1} that no such tower exists. Thus $d(\Sigma^{-2(p-1), -(p-1)}x) = v_0^{m-k}y$. Since  the map $d$ is $\Ext^{s,f,w}_{\mathcal{E}(1)^\vee_p}\left(\mM_p^F, \mM_p^F\right) \cong \mM_{p}^{F}[v_{0},v_{1}]$-linear, we get 
\[
d(cv_{0}^{n}\Sigma^{-2(p-1), -(p-1)}x_0 ) = cv_{0}^{m - k + n}y
\]
for all $c \in \mM_{p}^{F}$ and $n \ge 0$. It is clear from the stem degrees of the generators that no other differentials can occur. 

To finish the proof, relabel $\Sigma^{-2(p-1), -(p-1)}v_1x$ as $x$ and $\Sigma^{-2(p-1), -(p-1)}y$ as $y_{k-m+1}$. Finally, note that the extensions $v_{1}y_{i} = v_{0}y_{i+1}$ and $v_{1}\tau^{2}y_{i} = v_{0}\tau^{2}y_{i+1}$ must occur for exactly the same reasons as in the classical topological case (that is, by looking at the representatives of $v_1y_i$ and $v_0y_{i+1}$ in the chain complex computing $\Ext_{\cE(1)_p^\vee}^{*,1,*}\left(L_p(k),L_p(m)\right)$).
\end{proof}

\end{document}
